\begin{trapbox}
\paragraph{误区：把导数为零当成极值的充要条件}
\begin{center}
\begin{tikzpicture}
\begin{axis}[at={(0,0)},anchor=south west,width=6.1cm,height=5.6cm,
 axis lines=middle,xmin=-1.5,xmax=1.5,ymin=-2,ymax=2,
 ticks=none,xlabel={$x$},ylabel={$y$},title={$x^3$：驻点无极值},title style={font=\small}]
\addplot[blue,very thick,domain=-1.2:1.2,samples=80]{x^3};
\addplot[orange!85!black,dashed,domain=-0.7:0.7]{0};
\node[font=\small,anchor=east] at (axis cs:-0.2,1.2) {$f'(0)=0$};
\end{axis}
\begin{axis}[at={(6.2cm,0)},anchor=south west,width=6.1cm,height=5.6cm,
 axis lines=middle,xmin=-1.5,xmax=1.5,ymin=-2,ymax=2,
 ticks=none,xlabel={$x$},ylabel={$y$},title={$|x|$：尖点有极小值},title style={font=\small}]
\addplot[blue,very thick,domain=-1.3:1.3,samples=101]{abs(x)};
\addplot[only marks] coordinates {(0,0)};
\node[font=\small,anchor=west] at (axis cs:0.15,-1.2) {$f'(0)$ 不存在};
\end{axis}
\end{tikzpicture}
\end{center}

左图 $x^3$ 在 $0$ 处导数为零，但两侧仍递增，没有极值；右图 $|x|$ 在 $0$ 处不可导，却有严格极小值。
正确做法：把驻点和不可导点列作候选，再用定义或两侧变化判断。

\newpage
\paragraph{误区：把闭区间端点当成内部极值点}
\begin{center}
\begin{tikzpicture}
\begin{axis}[at={(0,0)},anchor=south west,width=6.1cm,height=5.6cm,
 axis lines=middle,xmin=-1.3,xmax=1.3,ymin=-2,ymax=1.5,
 xtick={1},ytick={1},xlabel={$x$},ylabel={$y$},title={端点最值：$y=x$},title style={font=\small}]
\addplot[blue,very thick,domain=0:1]{x};
\addplot[only marks] coordinates {(0,0) (1,1)};
\node[font=\small,anchor=west] at (axis cs:0.2,-0.8) {$x\in[0,1]$};
\end{axis}
\begin{axis}[at={(6.2cm,0)},anchor=south west,width=6.1cm,height=5.6cm,
 axis lines=middle,xmin=-1.3,xmax=1.3,ymin=-2,ymax=1.5,
 xtick=\empty,ytick={-1},xlabel={$x$},ylabel={$y$},title={断点：峰顶被挖走},title style={font=\small}]
\addplot[blue,very thick,domain=-1.1:1.1,samples=80]{-x^2};
\addplot[only marks,mark=*,mark options={fill=white,draw=blue},mark size=2.5pt] coordinates {(0,0)};
\addplot[only marks,mark=*,mark size=2.5pt] coordinates {(0,-1)};
\node[font=\small,anchor=west] at (axis cs:0.15,-1.65) {$f(0)=-1$};
\end{axis}
\end{tikzpicture}
\end{center}

本条采用内点极值的约定。$f(x)=x$ 在 $[0,1]$ 上，端点 $0,1$ 分别是最小值点和最大值点，却不是本条定义下的极值点。
求闭区间最值必须把端点加入候选，而不能只解 $f'=0$。

\paragraph{误区：不检查极值点处的连续性}
若 $f(x)=-x^2$（$x\ne0$），而 $f(0)=-1$，左右导数仍呈左正右负，原点却是严格极小值点。
因此用导数变号推出峰顶时，不能漏掉该点连续的前提；不连续时直接比较函数值。
\end{trapbox}
